\documentclass[12pt,a4paper]{article}

\usepackage{fontawesome5}
\providecommand{\faPencilAlt}{\faIcon{pencil-alt}}

\usepackage{styles/couleurs}
\usepackage{styles/mise_en_page}
\usepackage{styles/boites}
\usepackage{styles/maths}
\usepackage{styles/figures}
\usepackage{styles/evaluations}

\usepackage{enumitem}
\usepackage{array}
\usepackage{tabularx}
\usepackage{amsmath}
\usepackage{amssymb}

\setlist[itemize]{leftmargin=1.4em,itemsep=0.15em,topsep=0.2em}
\setlist[enumerate]{leftmargin=*,itemsep=0.25em,topsep=0.25em}

\begin{document}

\enteteinterro[Calculatrice interdite]{4e}{Divisibilité et nombres premiers}{15 min}{12}

\exo{3}[Représenter]
Complète avec « multiple » ou « diviseur ». \pts{3}

\begin{enumerate}[label=\alph*)]
  \item \(7\) est un \ldots{} de \(56\).
  \item \(72\) est un \ldots{} de \(8\).
  \item \(5\) est un \ldots{} de \(45\).
\end{enumerate}

\lignes{3}

\exo{3}[Calculer]
Sans poser de division, indique si \(4\,590\) est divisible par
\(2\), \(3\), \(5\), \(9\) et \(10\).
Justifie avec les critères. \pts{3}

\lignes{5}

\exo{4}[Raisonner]
Pour chaque entier, indique s'il est premier.
S'il ne l'est pas, donne un diviseur différent de \(1\)
et de lui-même. \pts{4}

\[
29\qquad 39\qquad 53\qquad 91
\]

\lignes{5}

\exo{2}[Communiquer]
Explique pourquoi \(1\) n'est pas un nombre premier
et pourquoi \(2\) est le seul nombre premier pair. \pts{2}

\lignes{3}

\end{document}
